% 20-minute presentation for the viva (17 slides), in the style of slides.tex. Vraj Shah's and
% Anshul D Jain's slides are rebuilt from their own decks. Compile from this folder with
% tectonic presentation.tex (or pdflatex). Tables are read from ../results/.
\documentclass[aspectratio=169,11pt]{beamer}
\usetheme{default}
\setbeamertemplate{navigation symbols}{}
\setbeamertemplate{footline}[frame number]
\usepackage{amsmath,amssymb,booktabs,graphicx}
% generated by experiments/report_tables.py from results/p5/summary/persistence.csv
\newcommand{\perslinear}{+0.10}
\newcommand{\persexponential}{0.00}

\newcommand{\by}[1]{\hfill{\normalfont\small\itshape #1}}

\title{Optimal Stopping with Memory}
\subtitle{Deep recurrent policies for American options under rough volatility\\
and multi-market backtesting}
\author{\small
  \begin{tabular}{llcll}
  Anamay Mittal & 2024ADPS0865G & \quad & Sachit Kapoor & 2024ADPS0858G\\
  Vraj Shah     & 2024ADPS0846G &       & Mridul Kwatra & 2024ADPS0819G\\
  Anshul D Jain & 2024ADPS0844G &       &               &
  \end{tabular}}
\institute{MAC F244 Stochastic Calculus --- Project Presentation\\BITS Pilani, K K Birla Goa Campus}
\date{15 November 2026}

\begin{document}

\begin{frame}
\titlepage
\end{frame}

% ------------------------------------------------------------------ 2
\begin{frame}{The question\by{Anamay Mittal}}
\begin{itemize}
\item \textbf{The problem.} Exercising an American option, or closing a trade, is an
  \textbf{optimal stopping} problem: each day, stop and collect the reward, or wait.
\item \textbf{The classical answer.} Work backwards from the last date, asking each day whether
  waiting is worth more than stopping (binomial tree, Longstaff--Schwartz, Deep Optimal
  Stopping). This assumes today's price is enough to decide (Markov).
\item \textbf{Why memory matters.} Real volatility is persistent and rough, so the path so far
  may carry information. Backward induction then needs a much larger state.
\item \textbf{OSPG} (Damera Venkata \& Bhattacharyya, NeurIPS 2023) learns the stopping rule
  directly, optionally with a recurrent network that reads the history.
\end{itemize}
\vspace{4pt}
\textbf{We test OSPG in three steps}
\begin{enumerate}
\item Option benchmarks with known exact prices.
\item Stopping problems with memory: fractional Brownian motion and rough Bergomi.
\item \textbf{Ten years of real market data} with transaction costs: NIFTY 50, S\&P 100, crypto.
\end{enumerate}
\end{frame}

% ------------------------------------------------------------------ 3
\begin{frame}{Optimal stopping: the problem and the optimal rule\by{Anamay Mittal}}
\begin{itemize}
\item \textbf{The problem.} Prices $S_0,\dots,S_H$; $\mathcal F_j$ is the information at date $j$.
  Stopping at $j$ pays $R_j$, discounted to date 0 (put: $R_j=e^{-rt_j}(K-S_j)^+$).
  A rule may use only the past, $\{\tau=j\}\in\mathcal F_j$. Goal: $V_0=\sup_\tau\mathbb E[R_\tau]$.
\item \textbf{The optimal rule.} Working backwards gives the Snell envelope, the value of acting
  optimally from date $j$ on:
  \[U_H=R_H,\qquad U_j=\max\big(R_j,\ C_j\big),\qquad C_j=\mathbb E[U_{j+1}\mid\mathcal F_j].\]
  Stop the first time stopping pays at least the \textbf{continuation value} $C_j$:
  $\tau^*=\min\{j:U_j=R_j\}$.
\item \textbf{Classical methods} estimate $C_j$ from today's state, one date at a time, backwards:
  Longstaff--Schwartz (regression), DOS (one network per date), FQI (one network for all dates).
\item With memory, $C_j$ depends on the whole path, and an error at one date carries into the next.
\end{itemize}
\end{frame}

% ------------------------------------------------------------------ 4
\begin{frame}{Exact prices: the free-boundary problem\by{Anamay Mittal}}
\begin{itemize}
\item In continuous time the American put solves
  \[\min\big(-\partial_tV-\mathcal LV,\ V-g\big)=0,\qquad
  \mathcal LV=rS\partial_SV+\tfrac12\sigma^2S^2\partial_{SS}V-rV.\]
\item Where waiting is optimal, the Black--Scholes equation holds; where exercising is, $V=g$,
  with payoff $g(S)=(K-S)^+$.
\item Our solver: Crank--Nicolson in $\log S$, with a penalty enforcing $V\ge g$.
\item \textbf{Agrees with a binomial tree to within $4\times10^{-4}$} on all twenty
  Longstaff--Schwartz test cases (e.g.\ $S_0=36$, $\sigma=0.2$, $T=1$: Bermudan $4.478$,
  American $4.487$).
\item These exact prices are what every simulation method is measured against.
\end{itemize}
\end{frame}

% ------------------------------------------------------------------ 5
\begin{frame}{When is holding optimal? Optional sampling\by{Anamay Mittal}}
\begin{itemize}
\item For any bounded stopping time $\tau\le H$:
  \begin{itemize}
  \item \textbf{martingale} (a fair game): $\mathbb E[M_\tau]=M_0$ --- no rule beats holding;
  \item \textbf{submartingale} (drifting upwards): $\mathbb E[M_\tau]\le\mathbb E[M_H]$ ---
    stopping at the deadline is optimal.
  \end{itemize}
\item So a clever exit can only pay if prices are not a fair or upward-drifting game.
\item This is the lens for the real-data backtest: we will see that nothing beats holding.
\end{itemize}
\end{frame}

% ------------------------------------------------------------------ 6
\begin{frame}{OSPG: learn the stopping rule directly\by{Anamay Mittal}}
\begin{itemize}
\item \textbf{Idea.} Instead of estimating $C_j$, learn the rule. A network outputs, at each date,
  $\varphi_j$ = probability of stopping, given it has not stopped yet ($\varphi_H=1$).
\item \textbf{Exact objective.} Stopping does not change the path, so
  \[\psi_j=\varphi_j\prod_{n<j}(1-\varphi_n),\qquad J(\theta)=\mathbb E\Big[\sum_{j}\psi_jR_j\Big].\]
\item \textbf{Training.} Policy gradient on pre-simulated paths:
  $\nabla_\theta J=\mathbb E\big[\sum_j(R_j-b)\,\psi_j\,\nabla_\theta\log\psi_j\big]$;
  the constant $b$ cuts noise without bias since $\sum_j\psi_j=1$.
  \textbf{No rollouts, no backward recursion.}
\item \textbf{Two policies.} OSPG-MLP sees today only ($2\times20$); OSPG-GRU reads the history
  ($1\times20$ GRU). After training: stop at the first date with $\varphi_j\ge\frac12$.
\item \textbf{Three fixes} not in the paper: a ``mostly wait'' starting bias, a small penalty on
  the raw output, a separate date-0 decision. Without them training collapses.
\end{itemize}
\end{frame}

% ------------------------------------------------------------------ 7
\begin{frame}{1. Replication: OSPG matches the exact prices\by{Anamay Mittal}}
\small
\textbf{Set-up.} Bermudan put (exact $4.478$), max-call on 20/50 assets (Becker et al.\ reference),
geometric call on 7/20 assets (exact). Out-of-sample values, three seeds.

\vspace{4pt}
\centering\scriptsize
% generated by experiments/summarize.py: method & mean (sd) per problem
\begin{tabular}{lrrrrr}
\toprule
Method & Put & Max-call $d{=}20$ & Max-call $d{=}50$ & Geo-call $d{=}7$ & Geo-call $d{=}20$ \\
\midrule
LSM & 4.458 (0.004) & 51.69 (0.20) & 69.83 (0.32) & 10.169 (0.012) & 9.899 (0.011) \\
DOS & 4.465 (0.007) & 51.69 (0.16) & 69.72 (0.29) & -- & -- \\
FQI-MLP & 4.467 (0.011) & 51.68 (0.12) & 69.67 (0.21) & 10.145 (0.068) & 9.874 (0.056) \\
FQI-GRU & -- & 51.73 (0.14) & 69.81 (0.37) & -- & -- \\
OSPG-MLP & 4.473 (0.008) & 51.80 (0.20) & 69.93 (0.28) & 10.255 (0.046) & 9.989 (0.017) \\
OSPG-GRU & 4.465 (0.015) & 51.72 (0.21) & 69.87 (0.25) & 10.237 (0.065) & 9.941 (0.030) \\
\midrule
Reference & 4.478 & 51.67 & 69.74 & 10.260 & 10.033 \\
Paper, DNN-OSPG & -- & 51.02 & 69.06 & 10.252 & 10.010 \\
\bottomrule
\end{tabular}


\vspace{4pt}\raggedright
\begin{itemize}\small
\item \textbf{Results.} OSPG-MLP is closest to the exact prices: put $-0.12\%$ (others $-0.24\%$ to
  $-0.45\%$); geometric calls $-0.05\%$, $-0.44\%$ (LSM, FQI $-0.9\%$ to $-1.6\%$). Max-call: all
  within $0.3\%$.
\item \textbf{Memory has a cost here.} OSPG-GRU trails OSPG-MLP by $0.1$--$0.5\%$ in 4 of 5
  problems ($z=2.2$--$3.7$): with nothing to remember, it must learn to ignore the history.
\end{itemize}
\end{frame}

% ------------------------------------------------------------------ 8
\begin{frame}{2. Stopping with memory: fBm and rough Bergomi\by{Mridul Kwatra}}
\small
\begin{itemize}
\item \textbf{fBm} $W^h$: increments anti-correlated for $h<\frac12$, positively correlated for
  $h>\frac12$; reward $R_j=W^h_{t_j}$.
\item \textbf{Rough Bergomi} American put: $v_t=\xi_0\exp(\eta\widetilde W^h_t-\frac12\eta^2t^{2h})$,
  $\widetilde W^h_t=\sqrt{2h}\int_0^t(t-s)^{h-1/2}dW_s$ (hybrid scheme).
\end{itemize}
\vspace{2pt}
\begin{columns}[T]
\begin{column}{.5\textwidth}\centering\scriptsize
fBm, mean test reward\\[2pt]
\resizebox{.95\linewidth}{!}{% generated by experiments/report_tables.py from results/fbm_medium_summary.csv
\begin{tabular}{rrrrr}
\toprule
$h$ & DOS & FQI-GRU & OSPG-MLP & OSPG-GRU \\
\midrule
0.05 & 0.700 & 1.297 & 1.153 & 1.286 \\
0.10 & 0.508 & 1.040 & 0.926 & 1.032 \\
0.30 & 0.187 & 0.366 & 0.334 & 0.362 \\
0.50 & -0.001 & -0.006 & -0.004 & -0.008 \\
0.70 & 0.165 & 0.202 & 0.191 & 0.193 \\
0.95 & 0.353 & 0.328 & 0.357 & 0.349 \\
\bottomrule
\end{tabular}
}
\end{column}
\begin{column}{.5\textwidth}\centering\scriptsize
Rough Bergomi put (pilot, se $0.31$--$0.37$)\\[2pt]
\resizebox{.95\linewidth}{!}{% generated by experiments/report_tables.py from results/e4_h_sweep_summary.csv
\begin{tabular}{rrrrr}
\toprule
$h$ & LSM & DOS & OSPG-MLP & OSPG-GRU \\
\midrule
0.07 & 7.935 & 7.985 & 8.136 & 7.932 \\
0.10 & 7.698 & 7.776 & 7.934 & 7.690 \\
0.30 & 7.362 & 7.435 & 7.507 & 7.367 \\
0.50 & 7.540 & 7.619 & 7.605 & 7.479 \\
\bottomrule
\end{tabular}
}
\end{column}
\end{columns}
\vspace{4pt}
\begin{itemize}
\item \textbf{fBm:} $\approx0$ at $h=\frac12$ (optional stopping); the recurrent policies are clearly
  ahead for $h\le0.3$. At $h=0.2$ we match the paper: OSPG-GRU $0.639$ (paper $0.64$), OSPG-MLP
  $0.578$ ($0.58$).
\item \textbf{Rough Bergomi pilot} (3{,}000 paths): no method distinguishable, no GRU advantage.
\end{itemize}
\end{frame}

% ------------------------------------------------------------------ 9
\begin{frame}{3. Real-data backtest: data and trade windows\by{Vraj Shah}}
\centering\small
\begin{tabular}{lll}
\toprule
Universe & Price source & History\\
\midrule
NIFTY 50 (50 stocks) & Yahoo Finance adjusted close & 2015--2024\\
S\&P 100 (100 of 101 tickers) & Yahoo Finance adjusted close & 2015--2024\\
Crypto (8 assets) & Binance spot daily close & Aug 2017--Dec 2024\\
\bottomrule
\end{tabular}

\vspace{8pt}\raggedright
\begin{itemize}\small
\item \textbf{Trade:} buy at a close; each day, sell or continue; exit by a 30-calendar-day deadline.
\item A new entry every five sessions, so windows overlap: \textbf{72{,}048 trade windows}.
\item \textbf{Data checks:} failures logged; gaps kept (no interpolation); incomplete windows skipped.
\item \textbf{No look-ahead:} at date $j$ a policy sees prices and rewards only up to $j$.
\end{itemize}
\end{frame}

% ------------------------------------------------------------------ 10
\begin{frame}{3. Real-data backtest: features, costs and rewards\by{Vraj Shah}}
\begin{columns}[T]
\begin{column}{.43\textwidth}
\small\textbf{Features at each decision}
\begin{itemize}\footnotesize
\item return since entry
\item running drawdown from the highest close
\item fraction of the horizon left
\item latest log return; lags of 1, 2, 5 days
\item volatility of the last 20 log returns
\item index and VIX, lagged one observation
\end{itemize}
\end{column}
\begin{column}{.55\textwidth}
\small\textbf{Costs per side (basis points)}\\[3pt]
\footnotesize
\begin{tabular}{lcccc}
\toprule
Market & Fees & Half-spread & Slippage & Total\\
\midrule
NIFTY 50 & 12 & 2 & 2 & 16\\
S\&P 100 & 1 & 1 & 1 & 3\\
Crypto & 10 & 2 & 2 & 14\\
\bottomrule
\end{tabular}

\vspace{4pt}
$r_j=\dfrac{S_j}{S_0}\,\dfrac{1-c_s}{1+c_b}-1$, \ at $0\times$, $1\times$, $2\times$ costs.
\end{column}
\end{columns}

\vspace{8pt}
\small\textbf{Two rewards.} Linear $R=r$, and exponential utility $U(r)=\big(1-e^{-ar}\big)/a$ with
$a=5$, fixed before testing: $+10\%\mapsto+0.079$ but $-10\%\mapsto-0.130$, so losses weigh more.
\end{frame}

% ------------------------------------------------------------------ 11
\begin{frame}{3. Real-data backtest: walk-forward folds and rules\by{Vraj Shah}}
\centering\small
\begin{tabular}{cccc}
\toprule
Fold & Training (expanding) & Validation & Test\\
\midrule
1 & months 1--36 & 37--42 & 43--48\\
2 & months 1--42 & 43--48 & 49--54\\
3 & months 1--48 & 49--54 & 55--60\\
\bottomrule
\end{tabular}

\vspace{6pt}\raggedright
\begin{itemize}\small
\item Train on the past, test on the future; each fold moves on 6 months. Windows that cross a
  boundary are purged. \textbf{31 folds} (12 NIFTY 50, 12 S\&P 100, 7 crypto); tests 2018--2024
  (crypto 2021--2024).
\item \textbf{Rules fixed in advance}, checked at daily closes: Hold to the deadline;
  take-profit/stop-loss at $+5\%/-3\%$; trailing stop at $5\%$ below the highest close.
\item \textbf{Learned methods:} Longstaff--Schwartz, OSPG-MLP, OSPG-GRU.
\end{itemize}
\end{frame}

% ------------------------------------------------------------------ 12
\begin{frame}{3. Real-data backtest: grid and inference\by{Anshul D Jain}}
\begin{itemize}
\item \textbf{The grid.} Six methods $\times$ two rewards $\times$ three cost levels $\times$ 31 folds
  $\times$ three seeds: \textbf{3{,}348 runs}, all completed.
\item \textbf{Paired comparisons.} Every method is tested on the same assets, dates and windows:
  36{,}637 test windows (NIFTY 50 10{,}137; S\&P 100 24{,}822; crypto 1{,}678).
\item \textbf{Inference.} Windows overlap and assets move together, so i.i.d.\ standard errors
  would be too small. Intervals come from a paired moving-block bootstrap over dates
  (30-day blocks, 2{,}000 replicates), with Holm correction within each market, reward and cost.
\item \textbf{Robustness.} With 15- or 60-day blocks, every comparison with Hold keeps its verdict.
\end{itemize}
\end{frame}

% ------------------------------------------------------------------ 13
\begin{frame}{3. Result: nothing beats holding to the deadline\by{Anshul D Jain}}
\centering\scriptsize
\resizebox{\linewidth}{!}{% generated by experiments/report_tables.py from results/p5/summary/market_cost_tables.csv
% base cost (1x): mean economic net return per window, percent
\begin{tabular}{llrrrrrrr}
\toprule
Market & Objective & Hold & TP/SL & Trailing & LSM & OSPG-MLP & OSPG-GRU & Windows \\
\midrule
NIFTY 50 & linear & 1.87 & 0.57 & 1.12 & 0.53 & 1.26 & 1.17 & 10{,}137 \\
NIFTY 50 & exponential & 1.87 & 0.57 & 1.12 & 0.56 & 0.46 & 0.47 & 10{,}137 \\
S\&P 100 & linear & 1.30 & 0.45 & 0.84 & 0.58 & 0.67 & 0.75 & 24{,}822 \\
S\&P 100 & exponential & 1.30 & 0.45 & 0.84 & 0.52 & 0.14 & 0.41 & 24{,}822 \\
Crypto & linear & 3.57 & 0.21 & 1.30 & 1.04 & 3.48 & 3.31 & 1{,}678 \\
Crypto & exponential & 3.57 & 0.21 & 1.30 & -0.28 & 0.06 & 0.03 & 1{,}678 \\
\bottomrule
\end{tabular}
}

\vspace{6pt}\normalsize
\begin{itemize}\small
\item Mean net return per window after costs (percent), base costs.
\item \textbf{Hold has the highest mean in every market, under both rewards}: $1.87\%$ (NIFTY 50),
  $1.30\%$ (S\&P 100), $3.57\%$ (crypto).
\end{itemize}
\end{frame}

% ------------------------------------------------------------------ 14
\begin{frame}{3. How large is the gap?\by{Anshul D Jain}}
\centering\scriptsize
\resizebox{\linewidth}{!}{% generated by experiments/report_tables.py from results/p5/summary/paired_statistics.csv
% base cost (1x), economic net return per window in percent; OSPG-GRU minus comparator
\begin{tabular}{llccccc}
\toprule
Market & Objective & Hold & TP/SL & Trailing & LSM & OSPG-MLP \\
\midrule
NIFTY 50 & linear & -0.70$^{*}$ [-1.07, -0.29] & +0.60 [-0.35, +1.35] & +0.05 [-0.81, +0.76] & +0.64$^{*}$ [+0.21, +1.21] & -0.08 [-0.28, +0.22] \\
NIFTY 50 & exponential & -1.40$^{*}$ [-2.18, -0.64] & -0.10 [-0.75, +0.30] & -0.65 [-1.34, -0.12] & -0.08 [-0.54, +0.37] & +0.01 [-0.17, +0.23] \\
S\&P 100 & linear & -0.55$^{*}$ [-0.73, -0.33] & +0.30 [-0.36, +1.10] & -0.08 [-0.63, +0.60] & +0.18 [-0.23, +0.74] & +0.09 [-0.11, +0.33] \\
S\&P 100 & exponential & -0.90$^{*}$ [-1.49, -0.38] & -0.04 [-0.52, +0.50] & -0.43 [-0.93, +0.17] & -0.11 [-0.34, +0.15] & +0.27$^{*}$ [+0.00, +0.52] \\
Crypto & linear & -0.26 [-1.10, +0.40] & +3.10 [-3.81, +7.14] & +2.02 [-3.62, +5.48] & +2.28 [-3.14, +5.14] & -0.16 [-1.31, +0.49] \\
Crypto & exponential & -3.54 [-7.92, +3.77] & -0.18 [-0.66, +0.47] & -1.27 [-2.43, +0.76] & +0.31 [-0.01, +0.72] & -0.04 [-0.19, +0.25] \\
\bottomrule
\end{tabular}
}

\vspace{6pt}\normalsize
\begin{itemize}\small
\item OSPG-GRU minus each comparator, percentage points, 95\% date-block bootstrap interval;
  $^{*}$Holm-adjusted $p<0.05$.
\item \textbf{Equities: OSPG-GRU trails Hold by $0.55$--$1.40$ points per window}
  (Holm $p\le0.012$). Crypto: intervals too wide to decide.
\end{itemize}
\end{frame}

% ------------------------------------------------------------------ 15
\begin{frame}{3. Where did memory help?\by{Anshul D Jain}}
\begin{columns}[T]
\begin{column}{.54\textwidth}\centering
\scriptsize OSPG-GRU minus OSPG-MLP, points per window\\
exponential reward, 95\% date-block bootstrap interval\\
$^{*}$Holm-adjusted $p<0.05$\\[4pt]
\begin{tabular}{llr}
\toprule
Market & Costs & GRU $-$ MLP [95\% interval]\\
\midrule
NIFTY 50 & base & +0.01 [$-$0.17, +0.23]\\
S\&P 100 & base & +0.27$^{*}$ [+0.00, +0.52]\\
Crypto & base & $-$0.04 [$-$0.19, +0.25]\\
\midrule
S\&P 100 & zero & +0.20 [$-$0.07, +0.43]\\
S\&P 100 & double & +0.18 [$-$0.08, +0.38]\\
\bottomrule
\end{tabular}
\end{column}
\begin{column}{.42\textwidth}
\begin{itemize}\small
\item \textbf{Markets:} at base cost, only the S\&P 100 interval lies above zero, and only just.
  Under the linear reward, no market shows a difference.
\item \textbf{Costs:} the S\&P 100 lead is positive at all three cost levels but clear only at
  base cost.
\item \textbf{Volatility persistence:} rank correlation with the GRU advantage is
  \perslinear{} (linear) and \persexponential{} (exponential): weak at most.
\end{itemize}
\end{column}
\end{columns}
\end{frame}

% ------------------------------------------------------------------ 16
\begin{frame}{Conclusions\by{Sachit Kapoor}}
\begin{itemize}
\item \textbf{OSPG works without backward induction}: on Markovian benchmarks it matches exact
  prices better than Longstaff--Schwartz and FQI.
\item \textbf{Memory helps only where there is memory to use}: strongly for fBm with
  $h\le0.3$; not in the rough Bergomi pilot; not reliably on real markets.
\item \textbf{Learned exit timing does not beat holding} after costs; holding earned
  $1.30$--$3.57\%$ per window, a positive drift. Consistent with optional sampling: for a
  submartingale, early exits give up drift.
\item \textbf{Limitations:} prices are lower bounds (no dual bounds); rough-volatility runs are
  reduced pilots; current index members (survivorship bias); stylized close-to-close execution;
  submartingale property not tested directly.
\end{itemize}
\end{frame}

% ------------------------------------------------------------------ 17
\begin{frame}{References and data}
\footnotesize
\begin{itemize}
\item N. Damera Venkata, C. Bhattacharyya. Deep recurrent optimal stopping. \emph{NeurIPS}, 2023.
\item S. Becker, P. Cheridito, A. Jentzen. Deep optimal stopping. \emph{JMLR} 20(74), 2019.
\item F. A. Longstaff, E. S. Schwartz. Valuing American options by simulation.
  \emph{Review of Financial Studies} 14(1), 2001.
\item C. Bayer, P. Friz, J. Gatheral. Pricing under rough volatility. \emph{Quantitative Finance} 16(6), 2016.
\item J. Gatheral, T. Jaisson, M. Rosenbaum. Volatility is rough. \emph{Quantitative Finance} 18(6), 2018.
\item M. Bennedsen, A. Lunde, M. S. Pakkanen. Hybrid scheme for Brownian semistationary processes.
  \emph{Finance and Stochastics} 21(4), 2017.
\item P. A. Forsyth, K. R. Vetzal. Quadratic convergence for valuing American options using a
  penalty method. \emph{SIAM J. Sci. Comput.} 23(6), 2002.
\item \textbf{Data:} Yahoo Finance (via \texttt{yfinance}) for NIFTY 50 and S\&P 100 constituents;
  Binance public spot prices and CoinGecko rankings for crypto. Code: \texttt{github.com/falconmath27/DROS}.
\end{itemize}
\end{frame}

\end{document}
